\documentclass[10pt,a4paper]{amsart}
\usepackage[margin=19mm]{geometry}
\usepackage{amsmath,amssymb,mathtools}
\usepackage[scr=rsfs,cal=euler]{mathalfa}
\usepackage{xfrac}
\usepackage[unicode,hidelinks,bookmarks=false]{hyperref}
\newcommand{\ie}{i.e.\ }
\newcommand{\cf}{cf.\ }
\newcommand{\resp}{respectively\ }
\newcommand{\Subsection}{\subsection}
\providecommand{\nobreakdash}{\nobreak\hbox{-}}
\setlength{\emergencystretch}{2em}
\multlinegap0pt
\usepackage{bm}
\usepackage{bbm}
\usepackage{dsfont}
\usepackage{xspace}
\usepackage{cancel}
\usepackage{mathtools}
\usepackage{amsthm}
\makeatletter
\@ifundefined{theorem}{\newtheorem{theorem}{Theorem}}{}
\@ifundefined{proposition}{\newtheorem{proposition}[theorem]{Proposition}}{}
\makeatother
\title[Analysis for idempotent states on quantum permutation groups]{Analysis for idempotent states on quantum permutation groups}
\author{J.P. McCarthy}
\date{Source: arXiv:2301.13423v3 (CC BY 4.0)}
\begin{document}
\maketitle

\noindent\emph{Source and licence.} Analysis for idempotent states on quantum permutation groups. Version arXiv:2301.13423v3 (CC BY 4.0), distributed under the Creative Commons Attribution 4.0 International licence, as recorded in the arXiv licence metadata for this exact version and shown on the abstract page https://arxiv.org/abs/2301.13423v3. Full rights notice: licenses/arxiv-2301.13423-CC-BY-4.0.txt.\par\smallskip

\noindent\textit{Editorial scope.} Counted unit: the Proposition whose source label is \texttt{None} in the article (source0.tex lines 1223--1230, source label \texttt{None}), delivered together with the article's own complete original proof of it (source0.tex lines 1231--1284, one \texttt{proof} environment of 54 lines). No mathematics is written, repaired or re-derived: statement and proof are the article's own text line for line. Required context retained uncounted: none. Every definition, display and label of those lines is retained unchanged. Omitted from this package and not redistributed: 1--1222, 1285--1324. Nothing inside a retained excerpt is omitted or altered: every retained body is delivered exactly as the article's lines are written, so no omission receipt of any kind accompanies this package. Citations: the retained excerpts carry no citation command at all, so no bibliography entry of the article is transcribed and no reference is redistributed. Reference adaptation: none was needed and none was made; every reference inside the delivered unit resolves inside it, so no unresolved reference or citation is produced. Printed slips kept literal: no printed word, symbol, hypothesis or number of the counted statement or of its proof was altered. Layout and class: the article's own class is \texttt{amsart} with packages including amssymb, amsthm, adjustbox, dsfont, epsfig, xy, hyperref; the wrapper is the lane's pinned \texttt{amsart} editorial wrapper at 10pt on A4, which restates the theorem-like environments the retained text needs and adds the article's own macro definitions that the retained text needs (none), with extra wrapper packages (bm, bbm, dsfont, xspace, cancel, mathtools). The delivered numbering is the unit's own. Assets: no retained span contains a figure environment, a graphics-inclusion command, a PDF-page inclusion, a table or a TikZ picture; no image file is copied into this package, the package declares no asset dependency and no image file is redistributed with it. Licence: version arXiv:2301.13423v3 of this article is distributed under the Creative Commons Attribution 4.0 International licence, as recorded in the arXiv licence metadata for this exact version and shown on the abstract page https://arxiv.org/abs/2301.13423v3. A full rights notice is delivered at licenses/arxiv-2301.13423-CC-BY-4.0.txt. Characters: every retained excerpt is pure ASCII, so the delivered engine, XeLaTeX with its default Latin Modern text font, reports no missing character.
\begin{proposition}
  Let $\mathbb{G}$ and $\mathbb{H}$ be compact quantum groups and $\pi : C(\mathbb{G}) \to C_{\text{u}}(\mathbb{H})$ a surjective unital $*$-homomorphism satisfying $(\pi \otimes \pi) \circ \Delta = \Delta_{\mathbb{H}} \circ \pi$.  Denote by $h_{\mathbb{H}} = h_{C_{\text{u}}(\mathbb{H})} \circ \pi$. Let $\varphi$ be a state on $C(\mathbb{G})$. Then the following statements are equivalent.
\begin{enumerate}
\item[(i)] There exists a state $\psi$ on $C_{\text{u}}(\mathbb{H})$ such that $\varphi = \psi \circ \pi$.
\item[(ii)] $\varphi$ is in the quasi-subgroup generated by $h_{\mathbb{H}}$,
\item[(iii)] $\varphi \star h_{\mathbb{H}}=h_{\mathbb{H}}$.
\end{enumerate}
\end{proposition}

\begin{proof}
The implications $(i) \Rightarrow (ii) \Rightarrow (iii)$ are trivial. Assume that (iii) holds. Denote by $c_0(\widehat{\mathbb{G}})$ the $c_0$ direct sum of the matrix algebras $B(H_\alpha)$, $\alpha \in \operatorname{Irr}(\mathbb{G})$. Denote by $\ell^\infty(\widehat{\mathbb{G}})$ the $\ell^\infty$ direct sum. Denote by $W \in M(C(\mathbb{G}) \otimes c_0(\widehat{\mathbb{G}}))$ the corresponding direct sum of unitary representations of $\mathbb{G}$. Similarly define $V \in M(C(\mathbb{H}) \otimes c_0(\widehat{\mathbb{H}}))$ for the compact quantum group $\mathbb{H}$. The morphism $\pi$ dualizes to a morphism $\widehat{\pi} : \ell^\infty(\widehat{\mathbb{H}}) \to \ell^\infty(\widehat{\mathbb{G}})$ satisfying $(\pi \otimes \mathrm{id})(W) = (\mathrm{id} \otimes \widehat{\pi})(V)$.

\bigskip

If given, for $i = 1,2$, unitary representations $U_i \in M(C(\mathbb{G}) \otimes \mathcal{K}(H_i))$ of $\mathbb{G}$ on Hilbert spaces $H_i$, define the intertwiner space $\operatorname{Mor}_\mathbb{G}(U_1,U_2)$ as the space of bounded operators $A \in B(H_1,H_2)$ satisfying $(1 \otimes A) U_1 = U_2 (1 \otimes A)$. Every unitary representation $U \in M(C(\mathbb{G}) \otimes \mathcal{K}(H))$ of $\mathbb{G}$ can be restricted to a unitary representation $U_\mathbb{H} := (\pi \otimes \mathrm{id})(U)$ of $\mathbb{H}$. Denote by $\operatorname{Mor}_\mathbb{H}(U_1,U_2)$ the space of intertwiners between these restrictions, so that $A \in \operatorname{Mor}_\mathbb{H}(U_1,U_2)$ if and only if $A \in B(H_1,H_2)$ and $(1 \otimes A) U_{1,\mathbb{H}} = U_{2,\mathbb{H}} (1 \otimes A)$. Note that $\operatorname{Mor}_\mathbb{G}(U_1,U_2) \subset \operatorname{Mor}_\mathbb{H}(U_1,U_2)$. Write $\operatorname{End}_\mathbb{H}(U) = \operatorname{Mor}_\mathbb{H}(U,U)$.

\bigskip

Note that the formula $U = (\mathrm{id} \otimes \theta)(V)$ defines a bijective correspondence between unitary representations $U \in M(C(\mathbb{H}) \otimes \mathcal{K}(H))$ of $\mathbb{H}$ and unital normal $*$-homomorphisms $\theta : \ell^\infty(\widehat{\mathbb{H}}) \to B(H)$. Then $\operatorname{End}_\mathbb{H}(U) = \theta(\ell^\infty(\widehat{\mathbb{H}}))'$ and $\theta(\ell^\infty(\widehat{\mathbb{H}})) = \operatorname{End}_\mathbb{H}(U)'$. It follows that for $T = (T_\alpha)_{\alpha \in \operatorname{Irr}(\mathbb{G})}$ in $\ell^\infty(\widehat{\mathbb{G}})$,
\begin{equation}\label{eq.rel-com}
T \in \widehat{\pi}(\ell^\infty(\widehat{\mathbb{H}})) \quad\text{iff}\quad T_\alpha Z = Z T_\beta \;\;\text{for all $\alpha,\beta \in \operatorname{Irr}(\mathbb{G})$ and $Z \in \operatorname{Mor}_\mathbb{H}(\beta,\alpha)$.}
\end{equation}


Denote by $p_\varepsilon \in c_0(\widehat{\mathbb{H}})$ the minimal central projection that corresponds to the trivial representation of $\mathbb{H}$. Write $q = \widehat{\pi}(p_\varepsilon)$. Note that $(h_1 \otimes \mathrm{id})(W) = q$. Define $T \in \ell^\infty(\widehat{\mathbb{G}})$ by $T = (\varphi \otimes \mathrm{id})(W)$. Since $(\Delta \otimes \mathrm{id})(W) = W_{13} W_{23}$ and $\varphi \ast h_1 = h_1$, it follows that $q = T q$. Write $R = W(1 \otimes q) - (1 \otimes q)$. Because
$$R^* R = 2 (1 \otimes q) - (1 \otimes q) W (1 \otimes q) - (1 \otimes q) W^* (1 \otimes q) \; ,$$
it is the case that $(\varphi \otimes \mathrm{id})(R^* R) = 0$. Since $\varphi$ is a state on a C$^*$-algebra, it follows that
$$(\varphi \otimes \mathrm{id} \otimes \mathrm{id})(W_{12} R_{13}) = 0 \; .$$
This means that
\begin{equation}\label{eq.formula}
(\varphi \otimes \mathrm{id} \otimes \mathrm{id})(W_{12} W_{13}) \, (1 \otimes q) = T \otimes q \; .
\end{equation}
Write $T$ as the direct sum of the matrices $T_\alpha \in B(H_\alpha)$ for $\alpha \in \operatorname{Irr}(\mathbb{G})$. Similarly denote the components of $q$ by $q_\gamma$ for every $\gamma \in \operatorname{Irr}(\mathbb{G})$. Take arbitrary $\alpha,\beta,\gamma \in \operatorname{Irr}(\mathbb{G})$. Consider the component in $B(H_\beta) \otimes B(H_\gamma)$ of \eqref{eq.formula}. Since $W$ is the direct sum of the unitary representations $U_\zeta \in C(\mathbb{G}) \otimes B(H_\zeta)$, $\zeta \in \operatorname{Irr}(\mathbb{G})$,
\begin{equation}\label{eq.new-formula}
(\varphi \otimes \mathrm{id} \otimes \mathrm{id})(U_{\beta,12} U_{\gamma,13}) \, (1 \otimes q_\gamma) = T_\beta \otimes q_\gamma \; .
\end{equation}
Take arbitrary $\alpha \in \operatorname{Irr}(\mathbb{G})$, $X \in \operatorname{Mor}_\mathbb{G}(\alpha,\beta \otimes \gamma)$ and $Y \in \operatorname{Mor}_\mathbb{H}(\varepsilon,\gamma)$. Multiplying \eqref{eq.new-formula} on the left by $X^*$ and on the right by $1 \otimes Y$,
\begin{equation}\label{eq.new-new-formula}
(\varphi \otimes \mathrm{id} \otimes \mathrm{id})((1 \otimes X^*) U_{\beta,12} U_{\gamma,13}) \, (1 \otimes q_\gamma Y) = X^*(T_\beta \otimes q_\gamma Y) \; .
\end{equation}
Note that $q_\gamma Y = \widehat{\pi}(p_\varepsilon)_\gamma Y = Y \widehat{\pi}(p_\varepsilon)_\varepsilon = Y$. Therefore, the left hand side of \eqref{eq.new-new-formula} equals
\begin{align*}
(\varphi \otimes \mathrm{id} \otimes \mathrm{id})((1 \otimes X^*) U_{\beta,12} U_{\gamma,13}) \, (1 \otimes Y) &= (\varphi \otimes \mathrm{id})(U_\alpha (1 \otimes X^*)) \, (1 \otimes Y) \\ &= (\varphi \otimes \mathrm{id})(U_\alpha) \, X^* (1 \otimes Y) = T_\alpha \, X^* (1 \otimes Y) \; ,
\end{align*}
while the right hand side of \eqref{eq.new-new-formula} equals
$$X^* (T_\beta \otimes Y) = X^* (1 \otimes Y) \, T_\beta \; .$$
Thus
$$T_\alpha \, X^*(1 \otimes Y) = X^*(1 \otimes Y) \, T_\beta$$
for all $\alpha,\beta,\gamma \in \operatorname{Irr}(\mathbb{G})$ and for all $X \in \operatorname{Mor}_\mathbb{G}(\alpha,\beta \otimes \gamma)$ and $Y \in \operatorname{Mor}_\mathbb{H}(\varepsilon,\gamma)$. By taking linear combinations, the same holds if $\gamma$ is any finite dimensional unitary representation of $\mathbb{G}$, not necessarily irreducible.

\bigskip


Now $T_\alpha \, Z = Z \, T_\beta$ for all $\alpha,\beta \in \operatorname{Irr}(\mathbb{G})$ and $Z \in \operatorname{Mor}_\mathbb{H}(\beta,\alpha)$. To prove this, choose solutions of the conjugate equations for $\alpha,\beta \in \operatorname{Irr}(\mathbb{G})$, given by $t \in \operatorname{Mor}_\mathbb{G}(\varepsilon,\overline{\beta} \otimes \beta)$ and $s \in \operatorname{Mor}_\mathbb{G}(\varepsilon,\beta \otimes \overline{\beta})$. Writing $\gamma = \overline{\beta} \otimes \alpha$, $X = s \otimes 1$ and $Y = (1 \otimes Z)t$, it follows that $X^*(1 \otimes Y) = Z$ and the claim follows. By this claim and by \eqref{eq.rel-com}, this\ means that $T = \widehat{\pi}(S)$ for some $S \in \ell^\infty(\widehat{\mathbb{H}})$.

\bigskip

Every $a \in \mathcal{O}(\mathbb{G})$ is of the form $a = (\mathrm{id} \otimes \omega)(W)$ where $\omega$ is a uniquely determined, ``finitely supported'' functional on $\ell^\infty(\mathbb{G})$. Note that $\pi(a) = 0$ if and only if $\omega \circ \widehat{\pi} = 0$. Since $(\varphi \otimes \mathrm{id})(W) = T = \widehat{\pi}(S)$,  there is a well defined linear functional $\psi : \mathcal{O}(\mathbb{H}) \to \mathbb{C}$ such that $\psi(\pi(a)) = \varphi(a)$ for all $a \in \mathcal{O}(\mathbb{G})$.

\bigskip

Then $\psi(1) = \varphi(1) = 1$ and, because $\pi : \mathcal{O}(\mathbb{G}) \to \mathcal{O}(\mathbb{H})$ is surjective, $\psi(b^* b) \geq 0$ for all $b \in \mathcal{O}(\mathbb{H})$. Since $C(\mathbb{H})$ is the universal C$^*$-algebra of the compact quantum group $\mathbb{H}$, it follows that $\psi$ uniquely extends to a state on $C(\mathbb{H})$ (still denoted $\psi$). By construction, $\varphi = \psi \circ \pi$.
\end{proof}

\end{document}
